\documentclass[12pt, a4paper]{article}

% ─── Packages ────────────────────────────────────────────────────
\usepackage[utf8]{inputenc}
\usepackage[T1]{fontenc}
\usepackage{lmodern}
\usepackage[margin=1in]{geometry}
\usepackage{graphicx}
\usepackage{booktabs}
\usepackage{array}
\usepackage{tabularx}
\usepackage{amsmath, amssymb}
\usepackage{hyperref}
\usepackage{xcolor}
\usepackage{listings}
\usepackage{caption}
\usepackage{float}
\usepackage{enumitem}
\usepackage{titlesec}
\usepackage{fancyhdr}
\usepackage{setspace}
\usepackage{algorithm}
\usepackage{algpseudocode}
\usepackage{multirow}
\usepackage{parskip}

% ─── Colors ──────────────────────────────────────────────────────
\definecolor{codeblue}{RGB}{0, 102, 204}
\definecolor{codegray}{RGB}{100, 100, 100}
\definecolor{codegreen}{RGB}{34, 139, 34}
\definecolor{backcolour}{RGB}{245, 245, 250}
\definecolor{accentblue}{RGB}{41, 98, 255}

% ─── Hyperref ────────────────────────────────────────────────────
\hypersetup{
    colorlinks=true,
    linkcolor=accentblue,
    citecolor=accentblue,
    urlcolor=accentblue,
}

% ─── Code listings ───────────────────────────────────────────────
\lstdefinestyle{pythonstyle}{
    backgroundcolor=\color{backcolour},
    commentstyle=\color{codegreen},
    keywordstyle=\color{codeblue}\bfseries,
    basicstyle=\ttfamily\footnotesize,
    breaklines=true,
    captionpos=b,
    keepspaces=true,
    numbers=left,
    numbersep=5pt,
    numberstyle=\tiny\color{codegray},
    frame=single,
    rulecolor=\color{codegray!30},
    language=Python,
}
\lstset{style=pythonstyle}

% ─── Headers ─────────────────────────────────────────────────────
\pagestyle{fancy}
\fancyhf{}
\fancyhead[L]{\small RAG-VQA: Retrieval-Augmented Visual Question Answering}
\fancyhead[R]{\small Midterm Report}
\fancyfoot[C]{\thepage}
\renewcommand{\headrulewidth}{0.4pt}

% ─── Section formatting ─────────────────────────────────────────
\titleformat{\section}{\Large\bfseries\color{accentblue}}{\thesection}{1em}{}
\titleformat{\subsection}{\large\bfseries}{\thesubsection}{1em}{}
\titleformat{\subsubsection}{\normalsize\bfseries}{\thesubsubsection}{1em}{}

\setstretch{1.15}

% =====================================================================
\begin{document}

% ─── TITLE PAGE ──────────────────────────────────────────────────
\begin{titlepage}
    \centering
    \vspace*{1.5cm}
    {\Huge\bfseries RAG-VQA}\\[0.5cm]
    {\LARGE Retrieval-Augmented Visual Question Answering\\with BLIP-2}\\[2cm]
    {\Large\bfseries Computer Vision --- Midterm Project Report}\\[0.5cm]
    {\large Masters of Artificial Intelligence}\\[2.5cm]
    \begin{tabular}{c}
        \large\bfseries Team Members \\[0.3cm]
        \large Nandita Menon \\
        \large Jillian Priscilla \\
        \large Anushree Ashok \\
    \end{tabular}
    \vfill
    {\large March 2026}
\end{titlepage}

\tableofcontents
\newpage

% =====================================================================
\section{Abstract}
% =====================================================================

Visual Question Answering (VQA) models typically rely solely on the query image to produce an answer. When questions are ambiguous or require disambiguation, a single image may not provide sufficient signal. We augment \textbf{BLIP-2} (OPT-2.7B) with a \textbf{Retrieval-Augmented Generation (RAG)} pipeline that finds visually similar images from a knowledge base using \textbf{CLIP ViT-B/32} and \textbf{FAISS}, injects their Q\&A pairs as few-shot hints, and lets the model leverage this context for improved answers.

Our RAG pipeline boosts VQA accuracy from \textbf{42.33\%} to \textbf{51.33\%} (+9pp) on a 100-sample VQAv2 evaluation set. Error analysis reveals 63\% of remaining errors stem from the model \emph{ignoring} injected context, motivating \textbf{QLoRA fine-tuning} on a hints-style prompt format.

% =====================================================================
\section{Introduction}
% =====================================================================

\subsection{Problem Statement}

Standard VQA models encode image and question independently, then fuse them to predict an answer. This faces limitations when questions are ambiguous (``What sport is being played?'' with partial cues), when answer formatting varies, or when subtle visual details benefit from exemplar context.

\subsection{Proposed Solution}

\textbf{RAG-VQA} augments BLIP-2 with external visual evidence:
\begin{enumerate}[leftmargin=1.5cm]
    \item \textbf{Retrieve} visually similar images from a knowledge base using CLIP + FAISS.
    \item \textbf{Fuse} image and text similarity scores with weight $\alpha$.
    \item \textbf{Gate} the retrieval using confidence threshold $\tau$.
    \item \textbf{Inject} retrieved Q\&A pairs as few-shot context into BLIP-2's prompt.
    \item \textbf{Generate} a short answer conditioned on both query image and retrieved context.
\end{enumerate}

% =====================================================================
\section{Technology Stack}
% =====================================================================

\begin{table}[H]
\centering
\caption{Technology stack.}
\renewcommand{\arraystretch}{1.3}
\begin{tabularx}{\textwidth}{@{} l l X @{}}
\toprule
\textbf{Category} & \textbf{Technology} & \textbf{Purpose} \\
\midrule
VQA Model & BLIP-2 (OPT-2.7B) & Vision-language generator with Q-Former bridge \\
Retrieval & CLIP ViT-B/32 + FAISS & Dual-encoder embeddings (512-d) + cosine search \\
Fine-Tuning & QLoRA (4-bit NF4) + PEFT & Memory-efficient LoRA on q\_proj, v\_proj \\
Framework & PyTorch, HuggingFace & Model loading, tokenization, streaming datasets \\
Dataset & VQAv2 (lmms-lab) & COCO images, 10 human answers per question \\
Deployment & Netlify + Chart.js & Interactive project dashboard \\
\bottomrule
\end{tabularx}
\end{table}

% =====================================================================
\section{System Architecture}
% =====================================================================

\subsection{Pipeline Overview}

\begin{figure}[H]
\centering
\fbox{\parbox{0.95\textwidth}{
\centering
\textbf{RAG-VQA Pipeline}\\[0.4cm]
\begin{tabular}{ccccccccc}
\fbox{\parbox{1.5cm}{\centering\small\textbf{Input}\\Image +\\Question}}
& $\rightarrow$ &
\fbox{\parbox{2cm}{\centering\small\textbf{CLIP}\\Retriever\\(512-d)}}
& $\rightarrow$ &
\fbox{\parbox{2cm}{\centering\small\textbf{Fusion \&}\\Gating\\$\alpha$, $\tau$}}
& $\rightarrow$ &
\fbox{\parbox{2cm}{\centering\small\textbf{BLIP-2}\\Generator\\(4-bit)}}
& $\rightarrow$ &
\fbox{\parbox{1.5cm}{\centering\small\textbf{Output}\\1--3 word\\answer}}
\end{tabular}
}}
\caption{High-level pipeline architecture.}
\end{figure}

\subsection{Score Fusion \& Gating}

For each retrieval candidate, we compute a fused score:
\begin{equation}
    \text{score}_{\text{fused}} = \alpha \cdot s_{\text{img}} + (1 - \alpha) \cdot s_{\text{txt}}
\end{equation}

where $\alpha \in [0,1]$ is a manually swept hyperparameter. Candidates are ranked by fused score and top-3 are returned. A confidence gate checks if $\text{score}_{\text{fused}}^{(1)} \geq \tau$; if not, the system falls back to baseline inference without context.

\subsection{BLIP-2 Generation}

BLIP-2 has three components: (1) frozen ViT image encoder, (2) Q-Former that converts visual features into soft prompt tokens, and (3) frozen OPT-2.7B that generates the answer. We use two prompt formats:

\textbf{Few-shot} (pre-trained RAG):
\begin{lstlisting}[language={}]
Q: How many engines? A: 4
Q: How many planes? A: 1
Q: What airline does the plane fly for? A:
\end{lstlisting}

\textbf{Hints} (fine-tuned RAG):
\begin{lstlisting}[language={}]
Answer with 1-3 words. Hints may be wrong.
Hints:
- Similar image suggests: tennis
- Similar image suggests: racket
Question: What sport is being played?
Answer:
\end{lstlisting}

\subsection{4-Bit Quantization}

All inference and training uses NF4 quantization via BitsAndBytes, reducing memory from $\sim$10GB to $\sim$2.5GB and enabling single-GPU ($\leq$8GB VRAM) operation.

% =====================================================================
\section{Dataset and Evaluation}
% =====================================================================

We use \textbf{VQAv2} (lmms-lab/VQAv2) with COCO images and 10 human answers per question. To ensure zero data leakage, we partition the validation split:

\begin{table}[H]
\centering
\caption{Disjoint data splits.}
\renewcommand{\arraystretch}{1.3}
\begin{tabular}{@{} l l r l @{}}
\toprule
\textbf{Split} & \textbf{Range} & \textbf{Size} & \textbf{Purpose} \\
\midrule
Knowledge Base & val[0:5000] & 5,000 & CLIP-indexed retrieval corpus \\
Evaluation & val[5000:5100] & 100 & All ablation experiments \\
Fine-tune & val[5100:8100] & 3,000 & QLoRA hints-format training \\
\bottomrule
\end{tabular}
\end{table}

\textbf{Metric}: Official VQAv2 soft accuracy: $\text{acc} = \min(1,\; \text{matching\_answers} / 3)$, with normalization (lowercase, remove articles/punctuation).

% =====================================================================
\section{Experiments and Results}
% =====================================================================

\subsection{Ablation Study}

\begin{table}[H]
\centering
\caption{Ablation results (KB=5000, top\_k=3, candidate\_k=30).}
\renewcommand{\arraystretch}{1.3}
\begin{tabular}{@{} l c c c c c @{}}
\toprule
\textbf{Config} & $\boldsymbol{\alpha}$ & $\boldsymbol{\tau}$ & \textbf{RAG Used} & \textbf{Accuracy} & $\boldsymbol{\Delta}$ \\
\midrule
Baseline (no RAG) & --- & --- & 0\% & 42.33\% & --- \\
Text-only & 0.0 & 0.0 & 100\% & 41.33\% & \textcolor{red}{--1.00} \\
Balanced & 0.5 & 0.0 & 100\% & 46.33\% & \textcolor{green!60!black}{+4.00} \\
Balanced + gate & 0.5 & 0.5 & 71\% & 47.33\% & \textcolor{green!60!black}{+5.00} \\
\textbf{Image-only} & \textbf{1.0} & \textbf{0.0} & \textbf{100\%} & \textbf{51.33\%} & \textcolor{green!60!black}{\textbf{+9.00}} \\
Image-only + gate & 1.0 & 0.75 & 59\% & 49.33\% & \textcolor{green!60!black}{+7.00} \\
\bottomrule
\end{tabular}
\end{table}

\textbf{Key findings}: (1) Image-only retrieval ($\alpha{=}1.0$) achieves best accuracy; text-only \emph{hurts} performance. (2) All image-weighted configs outperform baseline (+4 to +9pp). (3) Gating trades peak accuracy for precision by filtering low-confidence retrievals.

\subsection{Error Analysis}

We categorize 46 incorrect predictions from the best config ($\alpha{=}1.0$, $\tau{=}0.0$):

\begin{table}[H]
\centering
\caption{Error breakdown (46 incorrect predictions).}
\renewcommand{\arraystretch}{1.3}
\begin{tabular}{@{} l r r p{7cm} @{}}
\toprule
\textbf{Category} & \textbf{Count} & \textbf{\%} & \textbf{Description} \\
\midrule
Model Ignored Context & 29 & 63\% & Relevant hints provided but OPT-2.7B has no learned mechanism to attend to few-shot Q\&A context. \\
Wrong Images Retrieved & 9 & 20\% & CLIP matched coarse scene features, missing fine-grained details the question targets. \\
Misleading Answers & 8 & 17\% & Visually similar images with Q\&A for a \emph{different} question, actively steering model wrong. \\
\bottomrule
\end{tabular}
\end{table}

The dominant failure (63\%) is the model \textbf{ignoring context}, not retrieval quality. This directly motivated QLoRA fine-tuning.

\subsection{QLoRA Fine-Tuning}

\begin{table}[H]
\centering
\caption{QLoRA configuration.}
\renewcommand{\arraystretch}{1.3}
\begin{tabular}{@{} l l @{}}
\toprule
\textbf{Parameter} & \textbf{Value} \\
\midrule
LoRA rank / alpha & $r{=}16$, $\alpha{=}32$ (scaling $2\times$) \\
Target modules & q\_proj, v\_proj in OPT-2.7B \\
Trainable params & 5.2M / 3.75B (\textbf{0.14\%}) \\
Quantization & 4-bit NF4, double quantization \\
Effective batch size & 2 $\times$ 8 grad\_accum = 16 \\
Learning rate & $1 \times 10^{-4}$, AdamW \\
Data & 3,000 samples; 15\% no-hints for baseline preservation \\
\bottomrule
\end{tabular}
\end{table}

\begin{table}[H]
\centering
\caption{Training progress.}
\renewcommand{\arraystretch}{1.3}
\begin{tabular}{@{} c c c c @{}}
\toprule
\textbf{Epoch} & \textbf{Avg Loss} & \textbf{Time} & \textbf{Sanity Exact-Match} \\
\midrule
1 & 3.9899 & 18.6 min & 0.00\% \\
2 & 2.9431 & 18.1 min & 1.00\% \\
3 & 2.7019 & 18.4 min & 3.00\% \\
\bottomrule
\end{tabular}
\end{table}

Loss drops from 3.99 to 2.70; the model transitions from completely ignoring hints to beginning to follow them, validating that QLoRA can teach context-awareness with only 0.14\% trainable parameters.

% =====================================================================
\section{Qualitative Examples}
% =====================================================================

\begin{table}[H]
\centering
\caption{RAG success cases ($\alpha{=}1.0$).}
\renewcommand{\arraystretch}{1.3}
\begin{tabular}{@{} p{4cm} l l l @{}}
\toprule
\textbf{Question} & \textbf{Baseline} & \textbf{RAG} & \textbf{Ground Truth} \\
\midrule
What airline does the plane fly for? & United Airlines & American Airlines & American Airlines \\
How are the eggs cooked? & not & scrambled & scrambled \\
Is it morning or evening? & morning & evening & evening \\
How much time is left? & None & 15 minutes & 15 minutes \\
Is the skier going downhill? & no & yes & yes \\
\bottomrule
\end{tabular}
\end{table}

% =====================================================================
\section{Reproducibility}
% =====================================================================

\begin{lstlisting}[language=bash, caption={Full reproduction pipeline.}]
pip install torch transformers peft bitsandbytes faiss-cpu datasets tqdm Pillow
python run.py infer --num_samples 100 --offset 5000
python run.py rag-build --num_samples 5000
python run.py rag-infer --num_samples 100 --alpha 1.0 --tau 0.0
python run.py eval
python run.py rag-make-train --num_samples 3000
python run.py rag-finetune --train_jsonl data/rag_train.jsonl --epochs 3
\end{lstlisting}

% =====================================================================
\section{Individual Contributions}
% =====================================================================

\subsection{Nandita Menon}

\begin{table}[H]
\centering
\renewcommand{\arraystretch}{1.3}
\begin{tabularx}{\textwidth}{@{} l X @{}}
\toprule
\textbf{Area} & \textbf{Contributions} \\
\midrule
RAG Pipeline & Designed and implemented the retrieval-augmented generation pipeline: CLIP dual-modal encoding, FAISS index construction, score fusion ($\alpha$-weighting), confidence gating ($\tau$ threshold), and the \texttt{RAGRetriever} class. \\
\addlinespace
QLoRA Fine-Tuning & Implemented the fine-tuning infrastructure: 4-bit NF4 quantization, LoRA adapter setup (q\_proj, v\_proj), training loop with gradient accumulation, checkpoint management. Designed the hints-style prompt format. \\
\addlinespace
BLIP-2 Integration & Implemented model loading with 4-bit quantization, baseline/few-shot/hints inference modes, and answer post-processing. \\
\addlinespace
Training Data & Built the training data generation pipeline (\texttt{rag\_make\_train\_jsonl.py}) with automated hint retrieval and 15\% no-hints ratio. \\
\addlinespace
Project Dashboard & Developed the Netlify dashboard with Chart.js visualizations, inference traces, demo gallery, and responsive design. \\
\addlinespace
CLI \& Integration & Designed the unified CLI (\texttt{run.py}) with subcommands for all pipeline stages. End-to-end integration and deployment. \\
\bottomrule
\end{tabularx}
\end{table}

\subsection{Jillian Priscilla}

\begin{table}[H]
\centering
\renewcommand{\arraystretch}{1.3}
\begin{tabularx}{\textwidth}{@{} l X @{}}
\toprule
\textbf{Area} & \textbf{Contributions} \\
\midrule
Ablation Experiments & Executed the ablation study across 6 configurations, varying $\alpha$ and $\tau$. Collected prediction outputs and debug context logs. \\
\addlinespace
Evaluation & Implemented VQA soft accuracy metric following official VQAv2 protocol, with per-question-type breakdown. \\
\addlinespace
Documentation & Prepared speaker notes for midterm presentation. Contributed to report writing. \\
\bottomrule
\end{tabularx}
\end{table}

\subsection{Anushree Ashok}

\begin{table}[H]
\centering
\renewcommand{\arraystretch}{1.3}
\begin{tabularx}{\textwidth}{@{} l X @{}}
\toprule
\textbf{Area} & \textbf{Contributions} \\
\midrule
Dataset Pipeline & Implemented VQAv2 streaming loader with offset-based disjoint splits ensuring zero data leakage. \\
\addlinespace
Error Analysis & Categorized 46 incorrect predictions into three failure modes (ignored context, wrong retrieval, misleading answers). \\
\addlinespace
Visualization & Created demo gallery image extraction, selected qualitative examples, and contributed to chart design. \\
\bottomrule
\end{tabularx}
\end{table}

\subsection{Contribution Summary}

\begin{table}[H]
\centering
\caption{Contribution matrix.}
\renewcommand{\arraystretch}{1.3}
\begin{tabular}{@{} l c c c @{}}
\toprule
\textbf{Component} & \textbf{Nandita} & \textbf{Jillian} & \textbf{Anushree} \\
\midrule
RAG Pipeline (retrieval, fusion, gating) & \textbf{Lead} & Support & Support \\
BLIP-2 Model \& Inference & \textbf{Lead} & Support & --- \\
QLoRA Fine-Tuning & \textbf{Lead} & Support & --- \\
Training Data Generation & \textbf{Lead} & --- & Support \\
Project Dashboard \& Deployment & \textbf{Lead} & Support & Support \\
CLI Integration (\texttt{run.py}) & \textbf{Lead} & --- & --- \\
Ablation Experiments & Support & \textbf{Lead} & Support \\
Evaluation Metric & Support & \textbf{Lead} & --- \\
Dataset \& Data Splits & --- & Support & \textbf{Lead} \\
Error Analysis & Support & --- & \textbf{Lead} \\
Visualization \& Demo Gallery & Support & --- & \textbf{Lead} \\
Report \& Presentation & Support & \textbf{Lead} & Support \\
\bottomrule
\end{tabular}
\end{table}

% =====================================================================
\section{Conclusion and Future Work}
% =====================================================================

\subsection{Conclusion}

\begin{enumerate}[leftmargin=1.5cm]
    \item RAG improves VQA accuracy by \textbf{+9pp} (42.33\% $\rightarrow$ 51.33\%) with image-only retrieval.
    \item Image similarity dominates; text-only retrieval \emph{hurts} performance.
    \item The bottleneck is context utilization (63\% ignored), not retrieval quality (20\%).
    \item QLoRA fine-tuning (0.14\% params) teaches context-awareness via hints prompts.
    \item Confidence gating ($\tau$) provides a precision--recall tradeoff.
\end{enumerate}

\subsection{Future Work}

\begin{itemize}[leftmargin=1.5cm]
    \item Scale KB to 50,000+ images; upgrade to CLIP ViT-L/14 (768-d).
    \item Replace manual $\alpha$ with a learned per-query fusion module.
    \item Extended fine-tuning with more epochs and full VQAv2 test evaluation.
    \item Larger LLM backbone (OPT-6.7B or FlanT5-XXL).
\end{itemize}

% =====================================================================
\section*{References}
\addcontentsline{toc}{section}{References}
% =====================================================================

\begin{enumerate}[label={[\arabic*]}, leftmargin=1.5cm]
    \item Li, J., et al. (2023). BLIP-2: Bootstrapping Language-Image Pre-training with Frozen Image Encoders and Large Language Models. \textit{ICML 2023}.
    \item Radford, A., et al. (2021). Learning Transferable Visual Models From Natural Language Supervision. \textit{ICML 2021}.
    \item Goyal, Y., et al. (2017). Making the V in VQA Matter. \textit{CVPR 2017}.
    \item Johnson, J., Douze, M., \& J\'{e}gou, H. (2019). Billion-scale similarity search with GPUs. \textit{IEEE Trans. Big Data}.
    \item Dettmers, T., et al. (2023). QLoRA: Efficient Finetuning of Quantized Language Models. \textit{NeurIPS 2023}.
    \item Hu, E. J., et al. (2022). LoRA: Low-Rank Adaptation of Large Language Models. \textit{ICLR 2022}.
    \item Lewis, P., et al. (2020). Retrieval-Augmented Generation for Knowledge-Intensive NLP Tasks. \textit{NeurIPS 2020}.
    \item Zhang, S., et al. (2022). OPT: Open Pre-trained Transformer Language Models. \textit{arXiv:2205.01068}.
\end{enumerate}

\end{document}
